\chapter{Indledning}

Iltmætningen i arterielt blod er en central fysiologisk parameter i anæstesiologi. Den kliniske guldstandard, $\text{SaO}_2$ målt ved blodgasanalyse, kræver en arteriel blodprøve og er dermed invasiv og uegnet til kontinuerlig overvågning. Pulsoximetri introducerer et ikke-invasivt alternativ ved at udnytte forskellen i lysabsorption mellem iltet og uiltet hæmoglobin ved typisk \SI{660}{nm} og \SI{940}{nm}. Forholdet mellem de transmitterede signaler omsættes via en empirisk kalibrering til et estimat $\text{SpO}_2$. Metoden er så billig og pålidelig, at den både er standard på sygehuse og for forbrugere \cite{Webster1997,Nitzan2014}.

Bag metoden ligger en række antagelser. Beer-Bouguer-Lamberts lov antager et homogent, ikke-spredende medium, mens biologisk væv både absorberer og spreder lys på en bølgelængdeafhængig måde. LED'er udsender desuden lys i et bånd, og den faktiske peakbølgelængde kan afvige fra databladets indskrevne værdi. Sammen med hudpigmentering, bevægelse, lav perfusion og dyshæmoglobiner gør disse forhold, at et pulsoximeter i praksis er resultatet af både bevidste designvalg og en empirisk kalibrering, som dækker over de antagelser, man ikke nemt kan løse analytisk \cite{Tuchin2015,Ruberry2025}.

Formålet med dette bachelorprojekt er at undersøge, hvordan et pulsoximetersystem kan opbygges fra grunden, hvordan bølgelængdevalget kan understøttes optisk, og hvilke fejlkilder der i praksis dominerer det endelige $\text{SpO}_2$-estimat. Projektet består af tre komplementære spor: 

\begin{enumerate}
    \item Et analogt pulsoximeter med multipleksede LED'er ved \SI{660}{nm} og \SI{940}{nm}, en Hamamatsu S1223 fotodiode og en analog signalkæde med TIA, Sallen-Key lavpasfilter og forstærkertrin.

    \item En hyperspektral opstilling med en bredbåndet Glo-bar og et lineært variabelt filter, der måler den relative vævstransmission ved udvalgte NIR-bølgelængder

    \item En spektrometrisk karakterisering af de anvendte lyskilder og et kommercielt pulsoximeter som reference.
\end{enumerate}